\documentclass[11pt]{article}
\usepackage[a4paper,margin=19mm]{geometry}
\usepackage[no-math]{fontspec}
\setmainfont{Latin Modern Roman}[SmallCapsFont=Latin Modern Roman Caps]
\usepackage{amsmath,amssymb,amsthm,mathtools,graphicx,hyperref}
\hypersetup{hidelinks}
\emergencystretch=3em
\usepackage{aliascnt,etoolbox}
\newcommand{\AicSharedCounter}[2]{\ifstrequal{#1}{proposition}{\ifstrequal{#2}{Proposition}{\newaliascnt{proposition}{theorem}\newtheorem{proposition}[proposition]{Proposition}\aliascntresetthe{proposition}}{\PackageError{aic2441-source-counter}{Unlisted environment or heading}{Only the exact eleven inventoried pairs are allowed}}}{\ifstrequal{#1}{lemma}{\ifstrequal{#2}{Lemma}{\newaliascnt{lemma}{theorem}\newtheorem{lemma}[lemma]{Lemma}\aliascntresetthe{lemma}}{\PackageError{aic2441-source-counter}{Unlisted environment or heading}{Only the exact eleven inventoried pairs are allowed}}}{\ifstrequal{#1}{claim}{\ifstrequal{#2}{Claim}{\newaliascnt{claim}{theorem}\newtheorem{claim}[claim]{Claim}\aliascntresetthe{claim}}{\PackageError{aic2441-source-counter}{Unlisted environment or heading}{Only the exact eleven inventoried pairs are allowed}}}{\ifstrequal{#1}{corollary}{\ifstrequal{#2}{Corollary}{\newaliascnt{corollary}{theorem}\newtheorem{corollary}[corollary]{Corollary}\aliascntresetthe{corollary}}{\PackageError{aic2441-source-counter}{Unlisted environment or heading}{Only the exact eleven inventoried pairs are allowed}}}{\ifstrequal{#1}{conjecture}{\ifstrequal{#2}{Conjecture}{\newaliascnt{conjecture}{theorem}\newtheorem{conjecture}[conjecture]{Conjecture}\aliascntresetthe{conjecture}}{\PackageError{aic2441-source-counter}{Unlisted environment or heading}{Only the exact eleven inventoried pairs are allowed}}}{\ifstrequal{#1}{fact}{\ifstrequal{#2}{Fact}{\newaliascnt{fact}{theorem}\newtheorem{fact}[fact]{Fact}\aliascntresetthe{fact}}{\PackageError{aic2441-source-counter}{Unlisted environment or heading}{Only the exact eleven inventoried pairs are allowed}}}{\ifstrequal{#1}{definition}{\ifstrequal{#2}{Definition}{\newaliascnt{definition}{theorem}\newtheorem{definition}[definition]{Definition}\aliascntresetthe{definition}}{\PackageError{aic2441-source-counter}{Unlisted environment or heading}{Only the exact eleven inventoried pairs are allowed}}}{\ifstrequal{#1}{problem}{\ifstrequal{#2}{Problem}{\newaliascnt{problem}{theorem}\newtheorem{problem}[problem]{Problem}\aliascntresetthe{problem}}{\PackageError{aic2441-source-counter}{Unlisted environment or heading}{Only the exact eleven inventoried pairs are allowed}}}{\ifstrequal{#1}{question}{\ifstrequal{#2}{Question}{\newaliascnt{question}{theorem}\newtheorem{question}[question]{Question}\aliascntresetthe{question}}{\PackageError{aic2441-source-counter}{Unlisted environment or heading}{Only the exact eleven inventoried pairs are allowed}}}{\ifstrequal{#1}{example}{\ifstrequal{#2}{Example}{\newaliascnt{example}{theorem}\newtheorem{example}[example]{Example}\aliascntresetthe{example}}{\PackageError{aic2441-source-counter}{Unlisted environment or heading}{Only the exact eleven inventoried pairs are allowed}}}{\ifstrequal{#1}{setup}{\ifstrequal{#2}{Setup}{\newaliascnt{setup}{theorem}\newtheorem{setup}[setup]{Setup}\aliascntresetthe{setup}}{\PackageError{aic2441-source-counter}{Unlisted environment or heading}{Only the exact eleven inventoried pairs are allowed}}}{\PackageError{aic2441-source-counter}{Unlisted environment or heading}{Only the exact eleven inventoried pairs are allowed}}}}}}}}}}}}}
\newfontfamily\AicFiniteOneFont{Noto Sans Math}
\newcounter{AicFiniteOneCalls}
\newcommand{\mathbbm}[1]{\ifstrequal{#1}{1}{\stepcounter{AicFiniteOneCalls}\mathord{\text{\upshape\AicFiniteOneFont\char"1D7D9}}}{\PackageError{aic2441-finite-one}{Unlisted indicator argument}{Only exact original literal1 is permitted}}}
\AtEndDocument{\ifnum\value{AicFiniteOneCalls}=4\relax\else\PackageError{aic2441-finite-one}{Indicator occurrence count mismatch}{Exactly four original indicator calls are required}\fi}
\usepackage{color}
\usepackage{amsmath, amssymb, amsthm}
\usepackage{mathrsfs}
\usepackage{mathtools}
\usepackage[noabbrev,capitalize,nameinlink]{cleveref}
\crefname{equation}{}{}
\usepackage[noadjust]{cite}
\usepackage{graphics}
\usepackage{pifont}
\usepackage{tikz}
% Original bbm import retained in archived author preamble; independent finite indicator renderer only.
%\usepackage[T1]{fontenc}
\DeclareMathOperator*{\argmax}{arg\,max}
\DeclareMathOperator*{\argmin}{arg\,min}
\usetikzlibrary{arrows.meta}
\newcommand{\cmark}{\text{\ding{51}}}
\newcommand{\xmark}{\text{\ding{55}}}
\usepackage{environ}
\usepackage{framed}
\usepackage{url}
\usepackage[labelfont=bf]{caption}
\usepackage{cite}
\usepackage{graphicx}
\usepackage{tcolorbox}
\usepackage{enumerate}
\allowdisplaybreaks[1]


% ------   Theorem Styles -------
\numberwithin{equation}{section}
\newtheorem{theorem}{Theorem}[section]
\AicSharedCounter{proposition}{Proposition}
\AicSharedCounter{lemma}{Lemma}
\AicSharedCounter{claim}{Claim}
\crefname{claim}{Claim}{Claims}

\AicSharedCounter{corollary}{Corollary}
\AicSharedCounter{conjecture}{Conjecture}
\newtheorem*{question*}{Question}
\AicSharedCounter{fact}{Fact}

\theoremstyle{definition}
\AicSharedCounter{definition}{Definition}
\AicSharedCounter{problem}{Problem}
\AicSharedCounter{question}{Question}
\newtheorem*{definition*}{Definition}
\AicSharedCounter{example}{Example}
\AicSharedCounter{setup}{Setup}
\theoremstyle{remark}
\newtheorem*{remark}{Remark}
\newtheorem*{notation}{Notation}

% ----- Delimiters ----
\newcommand{\norm}[1]{\bigg\lVert#1\bigg\rVert}
\newcommand{\snorm}[1]{\lVert#1\rVert}
\newcommand{\sang}[1]{\langle #1 \rangle}

\newcommand{\mb}{\mathbb}
\newcommand{\mbf}{\mathbf}
\newcommand{\mbm}{\mathbbm}
\newcommand{\mc}{\mathcal}
\newcommand{\mf}{\mathfrak}
\newcommand{\mr}{\mathrm}
\newcommand{\ol}{\overline}
\newcommand{\on}{\operatorname}
\newcommand{\wh}{\widehat}
\newcommand{\wt}{\widetilde}
\newcommand{\msc}{\mathscr}
\begin{document}
\begin{center}{\Large Stable item 2441: Subset-sum range bound from a large Bernoulli atom, with square-root epsilon exponent}\end{center}
\paragraph{Source and attribution.} Vishesh Jain, Ashwin Sah, and Mehtaab Sawhney, \emph{Anticoncentration versus the Number of Subset Sums}. Advances in Combinatorics 2021:6; exact publisher-linked arXiv:2101.07726v2 (14 June 2021); published 28 June 2021; DOI 10.19086/aic.24872; exact official native archive and matching published PDF retrieved 2 October 2026. \href{https://www.advancesincombinatorics.com/article/24872-anticoncentration-versus-the-number-of-subset-sums}{Publisher article}. Exact native version: \href{https://creativecommons.org/licenses/by/4.0/}{CC BY4.0}. The byte-identical published PDF separately prints/links \href{https://creativecommons.org/licenses/by/3.0/}{CC BY3.0}.
\paragraph{Selected original conclusion.} Select only the full original Theorem 1.2: for every \(\epsilon>0\), dimension \(n\ge\epsilon^{-1/2}\), and real vector \(\vec w\in\mathbb R^n\) with the original Bernoulli atom \(\rho(\vec w)\ge\exp(-\epsilon n)\), the subset-sum range satisfies \(|\mathcal R(\vec w)|\le\exp(C\sqrt\epsilon n)\) for an absolute constant. Retain the entire original theorem, proof, all boundary-coordinate cases, Gaussian moment and shift-ratio/tail-integration core. The original abstract/context uses base-two exponential notation, while the selected theorem uses natural exponentials; both conventions remain as authored. The bin-packing consequence and supporting lemmas are unselected or supporting context, never extra units.
\paragraph{Literal source correspondence.} Three literal source interfaces are retained visibly. In the projected-family proof, native 305 and 309--310 use \(n-i\) remaining coordinates, while the next estimate at 313 literally prints the factor \(n/2\); the complete adjacent inequalities and final \(O(\cdot)n\) bound remain unchanged. Native 239 defines \(\operatorname{Bin}(k)\) using \(k\) trials, while native 320 chooses \(k=\epsilon^{-1/2}/2\) without an explicit integer-rounding convention; no editorial rounding bridge is supplied. Native 343 prints \(32sy^2/k^2\), native 353 defines \(f(y)\) with the literal linear term \(32sy/k^2\), and the following Gaussian evaluation at 361 prints the quadratic \(w^2\) term. All three displays and the complete intervening coupling/Jensen chain are unchanged. The earlier and later definitions of \(z\), and the explicit empty-supremum convention at 306, stay literal. No constant, exponent, variable or iteration correction, strengthened claim, or generated argument is supplied.
\paragraph{Difficulty (editorial): H3.} The selected main bound requires the complete new boundary-coordinate fiber-counting construction, including coordinates equal to k+1 and the projected support families, beyond the precisely cited injectivity observation. Its quantitative closure also requires the full Gaussian-sign coupling and Jensen moment calculation together with the new shift-ratio tail estimate and moment integration. Those interacting nonroutine constructions remain the main theorem’s own research obligations, rather than an application of a black-box inverse Littlewood–Offord theorem or a separately counted supporting lemma.
\paragraph{Editorial source boundary.} All retained original mathematical prose, formulas, proof boundaries, labels, references, author definitions and the complete six-entry bibliography are editable and unchanged. The independent article/AMS/19mm/Latin Modern Roman and Caps wrapper replaces publisher front matter, with the original section-reset theorem and equation numbering. Genuine installed aliascnt supplies only the eleven inventoried shared-counter heading pairs, rejecting all other pairs. Four exact original indicator-one calls use the pinned installed Noto Sans Math double-struck digit-one glyph; glyph shape and spacing differ from the original bbm font, but the literal body tokens and subscripts remain. Only the exact argument 1 and exactly four occurrences are allowed. This is an independent finite typography renderer; no denied package acquisition or fake bbm package is used. Original classes, logos, full raw sources and preambles stay privately preserved and unexecuted. No mathematical proof, repair or independent refereeing is added. Native CC4 and published PDF CC3 are separate grants.
\clearpage
\section{Introduction} \label{sec:introduction}
For $\vec{w} := (w_1,\dots, w_n) \in \mb{R}^{n}$ and a real random variable $\xi$, recall that the L\'evy concentration function of $\vec{w}$ with respect to $\xi$ is defined for all $r\ge 0$ by
\[\mc{L}_{\xi}(\vec{w},r) = \sup_{\tau \in \mb{R}}\mb{P}[|w_1 \xi_1 + \dots + w_n \xi_n - \tau| \le r],\]
where $\xi_1,\dots, \xi_n$ are i.i.d.\ copies of $\xi$. In combinatorial settings (where $\vec{w} \in \mb{Z}^{n}$) a particularly natural and interesting case is when $r = 0$ and $\xi$ is a Bernoulli random variable, i.e., $\xi = 0$ with probability $1/2$ and $\xi = 1$ with probability $1/2$. For lightness of notation, we will denote this special case by
\[\rho(\vec{w}) = \mc{L}_{\on{Ber}(1/2)}(\vec{w},0) = \sup_{\tau\in\mb{R}}\mb{P}[\sang{\vec{w},\vec{\xi}}=\tau].\]
In this note, we study the following question.
\begin{question}
\label{question}
For a vector $\vec{w} = (w_1,\dots, w_n) \in \mb{R}^{n}$ with $\rho(\vec{w}) \ge \rho$, how large can the range
\[\mc{R}(\vec{w}) = \{w_1\xi_1 + \dots + w_n \xi_n: \xi_i \in \{0, 1\}\}\]
be?
\end{question}
The two extremal examples here are $\vec{w} = (0,0,\dots,0)$, which corresponds to $\rho(\vec{w}) = 1$, $|\mc{R}(\vec{w})| = 1$ and $\vec{w} = (1,10,\dots,10^{n-1})$, which corresponds to $\rho(\vec{w}) = 2^{-n}$, $|\mc{R}(\vec{w})| = 2^{n}$. Motivated by these examples, one may ask if there is a smooth trade-off between $\rho(\vec{w})$ and $|\mc{R}(\vec{w})|$. This turns out not to be the case. Indeed, for any $\epsilon > 0$, Wiman \cite{W17} gave an example of a $\vec{w} \in \mb{Z}^{n}$ for which $|\mc{R}(\vec{w})| \ge 2^{(1-\epsilon)n}$ and $\rho(\vec{w}) \ge 2^{-0.7447n}$. At the other end of the spectrum, when $\rho(\vec{w}) \ge 2^{-\epsilon n}$, the so-called inverse Littlewood--Offord theory \cite{RV08, TV09, NV11} \emph{heuristically} suggests that $\vec{w}$ is essentially contained in a low-rank generalized arithmetic progression of `small' volume so that $|\mc{R}(\vec{w})|$ is also `small'. However, the number of `exceptional elements' in the inverse Littlewood--Offord theorems \cite{RV08, TV09, NV11} is unfortunately too large to be able to rigorously establish such a statement.   

Nevertheless, in a recent work on the parameterized complexity of the bin packing problem (see \cref{sub:application}), Nederlof, Pawlewics, Swennenhuis and W{\k{e}}grzycki \cite{NPSW20} showed that for any $\epsilon > 0$, 
\[\rho(\vec{w}) \ge 2^{-\epsilon n} \implies |\mc{R}(\vec{w})| \le 2^{\delta(\epsilon)n},\]
where
\begin{equation}
\label{eqn:delta-old}
\delta(\epsilon) = O\left(\frac{\log\log(\epsilon^{-1})}{\sqrt{\log(\epsilon^{-1})}}\right).
\end{equation}
In particular, $\delta(\epsilon) \to 0$ as $\epsilon \to 0$. Moreover, we must have $\delta(\epsilon) \ge (2-o(1))\epsilon$, as can be seen by considering
\[\vec{w} = (C_1,\dots,C_1, C_2,\dots, C_2,\dots, C_{n/k},\dots,C_{n/k}) \in \mb{R}^{n},\]
where each $C_i$ is repeated $k$ times, and $C_i$ is sufficiently small compared to $C_{i+1}$ for all $i$. Indeed, for such $\vec{w}$, we have $\rho(\vec{w}) = 2^{-(\frac{1}{2}+o_k(1))\frac{n}{k}\log_2{k}}$ while $|\mc{R}(\vec{w})| \le 2^{(1+o_k(1))\frac{n}{k}\log_2{k}}$.

We conjecture that this example is essentially the worst possible, so that $\delta(\epsilon) \le 2\epsilon$. We are able to show that
\begin{equation}
\label{eqn:delta-new}
\delta(\epsilon) = O(\sqrt{\epsilon}),
\end{equation}
thereby obtaining an exponential improvement over \cref{eqn:delta-old}. More precisely, 
\begin{theorem}\label{thm:main}
Let $\epsilon>0$. For any $n\ge\epsilon^{-1/2}$ and any $\vec{w}\in\mb{R}^n$ satisfying $\rho(\vec{w})\ge \exp(-\epsilon n)$, we have
\[|\mc{R}(\vec{w})|\le\exp(C_{\ref{thm:main}}\epsilon^{1/2}n),\]
where $C_{\ref{thm:main}}$ is an absolute constant. 
\end{theorem}
We prove this theorem in \cref{sec:proof}.

\subsection{Application to bin packing}
\label{sub:application}
The bin packing problem is a classic NP-complete problem whose decision version may be stated as follows: given $n$ items with weights $w_1,\dots, w_n \in [0,1]$ and $m$ bins, each of capacity $1$, is there a way to assign the items to the bins without violating the capacity constraints? Formally, is there a map $f:[n] \to [m]$ such that 
$\sum_{i \in f^{-1}(j)}w_i \le 1$ for all $j \in [m]$?

Bj\"orklund, Husfeldt, and Koivisto \cite{BHK09} provided an algorithm for solving bin packing in time $\tilde{O}(2^{n})$ where the tilde hides polynomial factors in $n$. It is natural to ask whether the base of the exponent may be improved at all i.e.\ is there a (possibly randomized) algorithm to solve bin packing in time $\tilde{O}(2^{(1-\epsilon)n})$ for some absolute constant $\epsilon > 0$?

In recent work, Nederlof, Pawlewics, Swennenhuis and W{\k{e}}grzycki \cite{NPSW20} showed that this is true provided that the number of bins $m$ is fixed. More precisely, they showed that there exists a function $\sigma: \mb{N} \to \mb{R}^{>0}$ and a randomized algorithm for solving bin packing which, on instances with $m$ bins, runs in time $\tilde{O}(2^{(1-\sigma(m))n})$, where $\tilde{O}$ hides polynomials in $n$ as well as exponential factors in $m$. Their analysis, which crucially relies on \cref{eqn:delta-old}, gives a very small value of $\sigma(m)$ satisfying
\begin{equation}
    \label{eqn:sigma-old}
    \sigma(m) \le 2^{-m^9}. 
\end{equation}
Using \cref{thm:main} instead of \cref{eqn:delta-old} in a black-box manner in the analysis of \cite{NPSW20}, we exponentially improve the bound on $\sigma(m)$.
\begin{corollary}
\label{cor:bin-packing}
With notation as above, the randomized algorithm of \cite{NPSW20} solves bin packing instances with $m$ bins in time $\tilde{O}(2^{(1-\sigma(m))n})$ with high probability, for $\sigma \colon \mb{N} \to \mb{R}^{> 0}$ satisfying
\begin{equation}
    \label{eqn:sigma-new}
    \sigma(m) = \tilde{\Omega}(m^{-12}),
\end{equation}
where $\tilde{\Omega}$ hides logarithmic factors in $m$.
\end{corollary}

\begin{remark}
This follows by noting that the function $f_C(m)$ in \cite[Section~3.6]{NPSW20} is $\tilde{\Theta}(m^{-2})$ so that $\delta$ in \cite[Section~3.6]{NPSW20} is $\tilde{\Theta}(m^{-3})$. With \cref{thm:main}, the function $\varepsilon(\delta)$ in the runtime analysis of \cite[Section~3.4]{NPSW20} satisfies $\varepsilon(\delta) = O(\delta^2)$. Therefore, the function $f_B(\delta)$ in the same section is $\tilde{O}(\delta^{4})$, which is $\tilde{\Omega}(m^{-12})$. Note that if one were able to establish the conjecturally optimal bound $\delta = O(\epsilon)$, this would lead to $f_B(\delta) = \tilde{O}(\delta^{2})$, thereby giving the quadratically better $\sigma(m) = \tilde{\Omega}(m^{-6})$.
\end{remark}

\subsection{Notation}\label{sub:notation}
We use big-$O$ notation to mean that an absolute multiplicative constant is being hidden. We use $\on{Ber}(1/2)$ to denote the balanced $\{0,1\}$ Bernoulli distribution and $\on{Bin}(k)$ to denote the binomial distribution on $k$ trials with parameter $1/2$. Recall that $\on{Bin}(k)$ is the sum of $k$ independent $\on{Ber}(1/2)$ random variables. Given a distribution $\mu$, we let $\mu^{\otimes n}$ denote the distribution of a random vector with $n$ independent samples from $\mu$ as its coordinates. We also use the following standard additive combinatorics notation: $C+D = \{c+d: c\in C,d\in D\}$ is the sumset (if $C,D$ are subsets of the same abelian group), and for a positive integer $k$, we let $k\cdot C = C+\cdots+C$ ($k$ times) be the iterated sumset. Finally, in some cases we will use the notation $\Sigma\cdot$ or $\int\cdot$ to denote that the expression in the sum or integral is the same as in the previous line to simplify the presentation of long expressions.

\subsection{Outline of the proof}
\label{sub:outline}
As in \cite{NPSW20}, the starting point of our proof is the following observation: let $A$ denote a fixed (but otherwise arbitrary) set of unique preimages for points in $\mc{R}(\vec{w})$ (hence, $|A| = |\mc{R}(\vec{w})|$) and let $B$ denote the the set of preimages of a value $\tau \in \mb{R}$ realising $\rho(\vec{w})$. Then (\cref{lem:injective}) for any $k\geq 1$, the map $A \times (k\cdot B) \to A + k\cdot B$ is a bijection. In particular, if $\vec{a}$ is sampled from the uniform distribution on $A$ and $\vec{b}_1,\dots, \vec{b}_k$ are independently sampled from the uniform distribution on $B$, then
\begin{align*}
    |A| &= |A| \cdot \mb{P}[\vec{a} + \vec{b}_1 + \dots  + \vec{b}_k \in \{0,\dots, k+1\}^{n}]\\
        &= |A| \cdot \sum_{\vec{x} \in \{0,\dots, k+1\}^{n}}\mb{P}[\vec{a} + \vec{b}_1 + \dots + \vec{b}_k = \vec{x}]\\
        &\leq |A| \cdot \sum_{\vec{x} \in \{0,\dots, k+1\}^{n}} \mb{P}[\vec{a} = \vec{a}(\vec{x})] \cdot \mb{P}[\vec{b}_1 + \dots + \vec{b}_k = \vec{x} - \vec{a}(\vec{x})]\\
        &\leq \sum_{\vec{x} \in \{0,\dots, k+1\}^{n}}\mb{P}[\vec{b}_1 + \dots + \vec{b}_k = \vec{x} - \vec{a}(\vec{x})] 
\end{align*}


In \cite{NPSW20}, the largeness of $B$ is exploited by finding, for every $a \in A$, a large subset of $B$ which is `balanced' (in a certain sense) with respect to $a$. Instead, we exploit the largeness of $B$ directly by using the observation that the density of the uniform measure on $B$ with respect to the uniform measure on $\{0,1\}^{n}$ is at most $2^{n}/|B| \leq 2^{\varepsilon n}$. In particular, if we let $\mu_{k}$ denote the measure on $k\cdot B$ induced by the product measure on $B \times \dots \times B$ via the map $(b_1,\dots, b_k) \mapsto b_1+\dots +b_k$ and if we let $\on{Bin}(k)^{\otimes n}$ denote the $n$-fold product of the $\on{Binomial}(k,1/2)$ distribution, then the density of $\mu_{k}$ with respect to $\on{Bin}(k)^{\otimes n}$ is at most $2^{k\varepsilon n}$. This allows us to replace the measure $\mu_{k}$ appearing in the last line of the above equation by $\on{Bin}(k)^{\otimes n}$, at the cost of a factor of $2^{k\varepsilon n}$. Thus,
\[|A| \leq 2^{k\epsilon n}\cdot \sum_{\vec{x} \in \{0,\dots,k+1\}^{n}}\mb{P}_{\vec{x} \sim \on{Bin}(k)^{\otimes n}}[\vec{x} - \vec{a}(\vec{x})]\]

The above expression is still complicated by the presence of the shift $\vec{a}(\vec{x})$, about which we have no information except that it lies in the set $A$. The key technical lemma in the proof is \cref{lem:sup-ratio}, which essentially allows us to remove this shift after paying a factor which depends on $|A|$. Ultimately, this gives an upper bound on the sum in terms of $|A|$ and $k$, which amounts to an upper bound on $|A|$ in terms of $k, \epsilon$, and $|A|$. Optimizing the value of the free parameter $k$ now gives the desired conclusion. 

\section{Proof of \texorpdfstring{\cref{thm:main}}{Theorem 1.2}}
\label{sec:proof}


We begin by recording the following key comparison bound, which will be proved at the end of this section.

\begin{lemma}\label{lem:sup-ratio}
Let $n\ge k \ge C_{\ref{lem:sup-ratio}}$, where $C_{\ref{lem:sup-ratio}}$ is a sufficiently large absolute constant and let $\delta > 0$. For any $A \subseteq \{0,1\}^{n}$ with $|A| \le \exp(\delta n)$, the following holds.  
Let $\vec{x},\vec{b}\sim\on{Bin}(k)^{\otimes n}$ be independent $n$-dimensional random vectors. Then, 
\[\mb{E}_{\vec{x}}\bigg[\sup_{\vec{a}\in A}\frac{\mb{P}_{\vec{b}}[\vec{b} = \vec{x}-\vec{a}]}{\mb{P}_{\vec{b}}[\vec{b} = \vec{x}]}\bigg]\le\exp\left(C_{\ref{lem:sup-ratio}}\left(\frac{1}{k}+\sqrt{\frac{\delta}{k}}\right)n\right).\]
\end{lemma}

Let $n$, $\epsilon$, and $\vec{w}$ be as in \cref{thm:main}. Let $\tau$ be such that $\mb{P}[\sang{\vec{w}, \vec{\xi}} = \tau] = \rho(\vec{w})$, where $\vec{\xi}$ is a random vector with i.i.d.\ $\on{Ber}(1/2)$ components. Let 
\[B = \{\vec{\xi} \in \{0,1\}^{n} : \sang{\vec{w}, \vec{\xi}} = \tau\}.\]
In particular, $|B|\ge \exp(-\epsilon n)\cdot 2^{n}$.
Let $|\mc{R}(\vec{w})| = \exp(\delta n)$. For each $r \in \mc{R}(\vec{w})$, let $\vec{\xi}(r)$ be a fixed (but otherwise arbitrary) element of $\{0,1\}^{n}$ such that $\sang{\vec{w},\vec{\xi}(r)} = r$. Let
\[A =\{\vec{\xi}(r) \in \{0,1\}^{n}: r \in \mc{R}(\vec{w})\}.\]
Note that, by definition, for any distinct $\vec{a}_1, \vec{a}_2 \in A$, we have that $\sang{\vec{w}, \vec{a}_1} \neq \sang{\vec{w}, \vec{a}_2}$ and that $|A| = |\mc{R}(\vec{w})| = \exp(\delta n)$.

% We are now ready to prove \cref{thm:main}. We first define the crucial objects of study. We fix $\vec{w}\in\mb{R}^n$ satisfying the conditions of the theorem, which implies the existence of $\tau,\alpha\in\mb{R}$, and $A,B\subseteq\{\pm 1\}^n$ satisfying the following:
% \begin{itemize}
%     \item $\sang{\vec{w},\vec{a}_1}\neq\sang{\vec{w},\vec{a}_2}$ for distinct $\vec{a}_1,\vec{a}_2\in A$,
%     \item $\sang{\vec{w},\vec{b}} = \tau$ for all $\vec{b}\in B$, and
%     \item $|B|/2^n\ge \exp(-\epsilon n)$.
% \end{itemize}
We will make use of the simple, but crucial, observation from \cite{NPSW20} that $A$ and $k\cdot B$ have a full sumset for all $k\ge 1$. 
\begin{lemma}[{\cite[Lemma~4.2]{NPSW20}}]\label{lem:injective}
The map $(\vec{a},\vec{c})\mapsto\vec{a}+\vec{c}$ from $A\times(k\cdot B)$ to $A + k\cdot B$ is injective.
\end{lemma}
\begin{proof}
Indeed, if $\vec{a}_{1} + (\vec{b}^{(1)}_1 + \dots + \vec{b}^{(1)}_k) = \vec{a}_{2} + (\vec{b}^{(2)}_1 + \dots + \vec{b}^{(2)}_k$), where $\vec{a}_i \in A$ and $\vec{b}^{(i)}_j \in B$, then taking the inner product of both sides with $\vec{w}$ and using $\sang{\vec{w}, \vec{b}} = \tau$ for all $b \in B$, we see that $\sang{\vec{w}, \vec{a}_1} = \sang{\vec{w}, \vec{a}_2}$, which implies that $\vec{a}_1 = \vec{a}_2$ by the definition of $A$. 
\end{proof}

We are now ready to prove \cref{thm:main}.

\begin{proof}[Proof of \cref{thm:main}]
Let $k\ge 2$ be a parameter which will be chosen later depending on $\epsilon$. We may assume $\epsilon\in(0,(2C_{\ref{lem:sup-ratio}})^{-2})$ by adjusting $C_{\ref{thm:main}}$ appropriately at the end to make larger values trivial. By \cref{lem:injective}, for each  $\vec{x}\in\{0,\ldots,k+1\}^n$ for which there exist $\vec{a} \in A$ and $\vec{c} \in k\cdot B$ with $\vec{a}+\vec{c} = \vec{x}$, there exists a unique such choice $\vec{a} = \vec{a}(\vec{x}) \in A$. (For $\vec{x}\notin A+k\cdot B$, we let $\vec{a}(\vec{x})$ be an arbitrary element of $A$.) 

Now, let $\vec{a}$ be uniform on $A$, let $\vec{b}_1,\ldots,\vec{b}_k$ be uniform on $B$, and let $\vec{v}_1,\ldots,\vec{v}_k$ be uniform on $\{0,1\}^n$. Let $C_i\subseteq\{0,\ldots,k+1\}^n$ be the set of vectors with $i$ coordinates equal to $k+1$. For $\vec{x}\in\{0,\ldots,k+1\}^n$, we let $\vec{x}^\ast \in \{0,\dots,k\}^{n}$ denote the vector obtained by setting every occurrence of $k+1$ in $\vec{x}$ to $k$. We have
\begin{align*}
1 &= \mb{P}[\vec{a}+\vec{b}_1+\cdots+\vec{b}_k\in\{0,\ldots,k+1\}^n]\\
&= \sum_{i=0}^n\sum_{\vec{x}\in C_i}\mb{P}[\vec{a}+\vec{b}_1+\cdots+\vec{b}_k = \vec{x}]\\
&\le \sum_{i=0}^n\sum_{\vec{x}\in C_i}\mb{P}[\vec{a} = \vec{a}(\vec{x})]\mb{P}[\vec{b}_1+\cdots+\vec{b}_k = \vec{x}-\vec{a}(\vec{x})]\\
&\le\frac{1}{|A|}\sum_{i=0}^n\sum_{\vec{x}\in C_i}\bigg(\frac{2^n}{|B|}\bigg)^k\mb{P}[\vec{v}_1+\cdots+\vec{v}_k = \vec{x}-\vec{a}(\vec{x})]\\
&\le\frac{e^{k\epsilon n}}{|A|}\sum_{i=0}^n\sum_{\vec{x}\in C_i}\mb{P}[\vec{v}_1+\cdots+\vec{v}_k = \vec{x}^\ast]\sup_{\vec{a}\in A}\frac{\mb{P}[\vec{v}_1+\cdots+\vec{v}_k = \vec{x}-\vec{a}]}{\mb{P}[\vec{v}_1+\cdots+\vec{v}_k = \vec{x}^\ast]}\\
&= \frac{e^{k\epsilon n}}{|A|}\sum_{i=0}^n(1/2^k)^i\sum_{S\in\binom{[n]}{i}}\mb{E}_{\vec{x}\sim\on{Bin}(k)^{\otimes ([n]\setminus S)}\times\{k+1\}^{S}}\bigg[\sup_{\vec{a}\in A}\frac{\mb{P}[\vec{v}_1+\cdots+\vec{v}_k = \vec{x}-\vec{a}]}{\mb{P}[\vec{v}_1+\cdots+\vec{v}_k = \vec{x}^\ast]}\bigg].
\end{align*}
Let $A_S$ be the set of elements in $A \subseteq \{0,1\}^{n}$ whose support contains $S$. Let
\[A'_S = \{\vec{a}' \in \{0,1\}^{[n]\setminus S}: \exists \vec{a} \in A_S \text{ with }\vec{a}|_{[n]\setminus S} = \vec{a}'\}.\]
Recall that $|A| = \exp(\delta n)$. Abusing notation so that the supremum of an empty set is $0$, we can continue the above chain of inequalities to get that
\begin{align*}
1&\le\frac{e^{k\epsilon n}}{|A|}\sum_{i=0}^n(1/2^k)^i\sum_{S\in\binom{[n]}{i}}\mb{E}_{\vec{x}\sim\on{Bin}(k)^{\otimes([n]\setminus S)}\times\{k+1\}^{S}}\bigg[\sup_{\vec{a}\in A}\frac{\mb{P}[\vec{v}_1+\cdots+\vec{v}_k = \vec{x}-\vec{a}]}{\mb{P}[\vec{v}_1+\cdots+\vec{v}_k = \vec{x}^\ast]}\bigg]\\
&= \frac{e^{k\epsilon n}}{|A|}\sum_{i=0}^n(1/2^k)^i\sum_{S\in\binom{[n]}{i}}\mb{E}_{\vec{x}\sim\on{Bin}(k)^{\otimes([n]\setminus S)}}\bigg[\sup_{\vec{a}\in A_S'}\frac{\mb{P}[\vec{v}_1+\cdots+\vec{v}_k = (\vec{x}-\vec{a})\times \{k\}^{S}]}{\mb{P}[\vec{v}_1+\cdots+\vec{v}_k = \vec{x} \times \{k\}^{S}]}\bigg]\\
&= \frac{e^{k\epsilon n}}{|A|}\sum_{i=0}^n(1/2^k)^i\sum_{S\in\binom{[n]}{i}}\mb{E}_{\vec{x}\sim\on{Bin}(k)^{\otimes([n]\setminus S)}}\bigg[\sup_{\vec{a}\in A_S'}\frac{\mb{P}[(\vec{v}_1+\cdots+\vec{v}_k)|_{[n]\setminus S} = \vec{x}-\vec{a}]}{\mb{P}[(\vec{v}_1+\cdots+\vec{v}_k)|_{[n]\setminus S} = \vec{x}]}\bigg]\\
&\le \frac{e^{k\epsilon n}}{|A|}\left(\sum_{i=0}^{n/2}\cdot + \sum_{i=n/2}^{n}2^{-ki}\cdot 2^{n}\cdot \left(\max_{\ell}\frac{\max\Big\{\binom{k}{\ell- 1},\binom{k}{\ell}\Big\}}{\binom{k}{\ell}}\right)^{n-i}\right)\\
&\le \frac{e^{k\epsilon n}}{|A|}\left(\sum_{i=0}^{n/2}\cdot + n\cdot 2^{-kn/2}\cdot 2^{n}\cdot k^{n}\right)\\
&\le\frac{e^{k\epsilon n}}{|A|}\bigg(\sum_{i=0}^{n/2}\binom{n}{i}2^{-ki}\exp(C_{\ref{lem:sup-ratio}}(k^{-1}+(2\delta)^{1/2}k^{-1/2})(n/2)) + 2^{-kn/4}\bigg)\\
&\le\exp(-\delta n)\exp\left(O(k\epsilon + k^{-1} + \delta^{1/2}k^{-1/2})n\right)
\end{align*}
by \cref{lem:sup-ratio} applied to $A_S$, as long as $n/2\ge k\ge C_{\ref{lem:sup-ratio}} \ge 20$. To deduce the last line, note that $\binom{n}{i}2^{-ki}\le (2^{-k}en/i)^i$, so for $i\ge\lceil en/2^{k-1}\rceil$ the sum of weighted binomials is bounded by a geometric series. Additionally, for $1\le i\le en/2^k$, if this interval is nonempty, the sum of binomials is certainly bounded by $\exp(O(k^{-1}n))$.

Hence, the above inequality yields
\[\delta\le C(k\epsilon + k^{-1}+ \delta^{1/2}k^{-1/2})\]
for some absolute constant $C > 0$. Now letting $k = \epsilon^{-1/2}/2$ (note that this satisfies $2C_{\ref{lem:sup-ratio}}\le 2k = \epsilon^{-1/2} \le n$), we find that
\[\delta = O(\epsilon^{1/2}),\]
 as desired.
\end{proof}

The proof of \cref{lem:sup-ratio} relies on the following preliminary estimate. 
\begin{lemma}\label{lem:initial-bound}
If $1\le s\le k/(16\pi)$, then \[\mb{E}_{x\sim\on{Bin}(k)}\bigg(\frac{x}{k+1-x}\bigg)^s\le\exp(10\pi s^2/k)+2k^s(4/5)^k.\]
\end{lemma}
\begin{proof}
We let $x\sim\on{Bin}(k)$ and $y = x-k/2\sim\on{Bin}(k)-k/2$ throughout. We let $z\sim\mc{N}(0,k\pi/8)$. 
%Note that there is a coupling of $y,z$ so that $y = z + z'$, where $\mb{E}[z'|z] = 0$. 
We have
\begin{align*}
\mb{E}_{x\sim\on{Bin}(k)}\bigg[\bigg(\frac{x}{k+1-x}\bigg)^s\bigg]&=\mb{E}_y\bigg[\bigg(1 + \frac{2y-1}{k/2+1-y}\bigg)^s\bigg]\\
&\le \mb{E}_y\bigg[\bigg(1+\frac{2y}{k/2+1-y}\bigg)^s\mbm{1}_{|y|\le k/3}\bigg] + k^s\mb{P}[|y|\ge k/3]\\
&\le\mb{E}_y\bigg[\bigg(1+\frac{2y}{k/2+1-y}\bigg)^s\mbm{1}_{|y|\le k/3}\bigg] + 2k^s(4/5)^k.
\end{align*}
Note that the probability estimate for $\mb{P}[\mbm{1}_{|y|\ge k/3}]$ follows from the sharp (entropy) version of the Chernoff-Hoeffding theorem. Since for $|y| \le k/3$,
\[\frac{2y}{(k/2 + 1-y)} \le \frac{2y}{k/2+1} + \frac{8y^{2}}{(k/2+1)^{2}},\]
and using $(1+x) \le \exp(x)$, we can continue the previous inequality as
\begin{align*}
\mb{E}_{x\sim\on{Bin}(k)}\bigg[\bigg(\frac{x}{k+1-x}\bigg)^s\bigg]&\le\mb{E}_y\bigg[\bigg(1+\frac{2y}{k/2+1}+\frac{8y^2}{(k/2+1)^2}\bigg)^s\mbm{1}_{|y|\le k/3}\bigg] + 2k^s(4/5)^k\\
&\le\mb{E}_y\bigg[\exp\bigg(\frac{4sy}{k+2}+\frac{32sy^2}{k^2}\bigg)\bigg] + 2k^s(4/5)^k.
\end{align*}
Now, let $z_1,\dots,z_k$ be i.i.d.\ $\mc{N}(0,1)$ random variables. Then,
\[y \sim \frac{1}{2}\left(\on{sgn}{z_1} + \dots + \on{sgn}{z_k}\right).\]
Moreover, for any $-k \le \ell \le k$,
\[\mb{E}[z_1 + \dots + z_k \mid \on{sgn}(z_1) + \dots + \on{sgn}(z_k) = \ell] = \sqrt{\frac{2}{\pi}}\ell.\]
In particular, under this coupling of $y,z_1,\dots,z_k$, we have
\[\mb{E}[z_1+\dots + z_k \mid y] = \sqrt{\frac{8}{\pi}}y.\]
Let $z = z_1 + \dots +z_k$, so that $z \sim \mc{N}(0,k)$.
Then, by the convexity of
\[f(y)= \exp\bigg(\frac{4sy}{k+2}+\frac{32sy}{k^2}\bigg)\]
and using Jensen's inequality, we have
\begin{align*}
\mb{E}_{y}f(y)
&= \mb{E}_{y, z_1, \dots, z_k}f(y)\\
&= \mb{E}_{y,z_1,\dots,z_k} f\left(\sqrt{\frac{\pi}{8}}\mb{E}[z \mid y]\right)\\
&\le \mb{E}_{z}f(\sqrt{\pi}z/\sqrt{8})\\
%=\mb{E}_z\exp\bigg(\frac{4sz}{k+2}+\frac{32sz^2}{k^2}\bigg)\\
&= \mb{E}_{w\sim\mc{N}(0,1)}\exp\bigg(\frac{s\sqrt{2k\pi}}{k+2}w+\frac{4s\pi}{k}w^2\bigg)\\
&= \bigg(1-\frac{8\pi s}{k}\bigg)^{-1/2}\exp\bigg(\frac{\pi s^2k^2}{(k+2)^2(k-8\pi s)}\bigg)\\
&\le\exp\bigg(\frac{8\pi s}{k} + \frac{2\pi s^2}{k}\bigg) \\
&\le\exp(10\pi s^2/k).\qedhere
\end{align*}
\end{proof}

Finally, we can prove \cref{lem:sup-ratio}
\begin{proof}[Proof of \cref{lem:sup-ratio}]
We may assume that $\delta \ge 2000/k$ since the statement for $\delta < 2000/k$ follows from the statement for $\delta = 2000/k$. Also, note that we may assume that $\delta \le \log2$.
For any $t \in \mb{R}$, we have
\begin{align*}
\mb{P}_{\vec{x}}\bigg[\sup_{\vec{a}\in A}\frac{\mb{P}_{\vec{b}}[\vec{b} = \vec{x}-\vec{a}]}{\mb{P}_{\vec{b}}[\vec{b} = \vec{x}]}\ge e^{tn}\bigg]&\le |A|\sup_{\vec{a}\in A}\mb{P}_{\vec{x}}\bigg[\frac{\mb{P}_{\vec{b}}[\vec{b} = \vec{x}-\vec{a}]}{\mb{P}_{\vec{b}}[\vec{b} = \vec{x}]}\ge e^{tn}\bigg]\\
&\le |A|\sup_{\vec{a}\in A}\inf_{s\ge 2}\exp(-stn)\mb{E}_{\vec{x}}\bigg[\bigg(\frac{\mb{P}_{\vec{b}}[\vec{b} = \vec{x}-\vec{a}]}{\mb{P}_{\vec{b}}[\vec{b} = \vec{x}]}\bigg)^s\bigg]\\
&= |A|\sup_{\vec{a}\in A}\inf_{s\ge 2}\exp(-stn)\prod_{i=1}^n\mb{E}_{x\sim\on{Bin}(k)}\bigg[\bigg(\frac{\mb{P}[\on{Bin}(k) = x-a_i]}{\mb{P}[\on{Bin}(k) = x]}\bigg)^s\bigg]\\
&\le |A|\inf_{s\ge 2}\exp(-stn)\bigg(\mb{E}_{x\sim\on{Bin}(k)}\bigg(\frac{x}{k+1-x}\bigg)^s\bigg)^n.
\end{align*}
In the last line, we have used that
\begin{align*}
    \mb{E}_{x\sim\on{Bin}(k)}\bigg[\bigg(\frac{x}{k+1-x}\bigg)^s\bigg]
    &\ge \bigg(\mb{E}_{x\sim\on{Bin}(k)}\bigg[\frac{x^2}{(k+1-x)^2}\bigg]\bigg)^{s/2}\\
    &= \bigg(\sum_{\ell = 0}^{k-1}\frac{\ell+1}{k-\ell}\binom{k}{\ell}2^{-k}\bigg)^{s/2}\\
    &= \bigg(\sum_{\ell = 0}^{k-1}\bigg(\frac{k+2}{k} + \frac{4(k+1)(\ell-k/2)}{k^2} + \frac{(k+1)(k-2\ell)^2}{k^2(k-\ell)}\bigg)\binom{k}{\ell}2^{-k}\bigg)^{s/2}\\
    &\ge \bigg(\sum_{\ell = 0}^{k-1}\bigg(\frac{k+2}{k} + \frac{4(k+1)(\ell-k/2)}{k^2}\bigg)\binom{k}{\ell}2^{-k}\bigg)^{s/2}\\
    &=\bigg(\frac{k+2}{k}-\frac{3k+4}{k}2^{-k}\bigg)^{s/2}\\
    &\ge 1\\
\end{align*}
if $k\ge 3$.
Therefore, by \cref{lem:initial-bound}, we have
\begin{align*}
\mb{P}_{\vec{x}}\bigg[\sup_{\vec{a}\in A}\frac{\mb{P}_{\vec{b}}[\vec{b} = \vec{x}-\vec{a}]}{\mb{P}_{\vec{b}}[\vec{b} = \vec{x}]}\ge e^{tn}\bigg]&\le |A|\inf_{s\ge 2}\exp(-stn)\bigg(\mb{E}_{x\sim\on{Bin}(k)}\bigg(\frac{x}{k+1-x}\bigg)^s\bigg)^n\\
&\le |A|\inf_{2\le s\le k/(16\pi)}\exp(-stn)\bigg(\exp(10\pi s^2/k)+2k^s(4/5)^k\bigg)^n\\
&\le |A|\inf_{2\le s\le k/(10\log k)}\exp(-stn)\bigg(\exp(12\pi s^2/k)\bigg)^n\\
& \le 
\begin{cases}
|A|\exp\left(-\frac{kt^{2}n}{48\pi}\right) &\quad \text{ if } \sqrt{\frac{96\pi\delta}{k}}\le t\le (\log{k})^{-1}\\
|A|\exp\left(-\frac{kn}{48\pi (\log{k})^{2}}\right) &\quad \text{ if } (\log{k})^{-1}\le t \le \log{k}.
\end{cases}
\end{align*}
Here, the second case follows by plugging in $s = k/(24\pi\log{k})$ and simplifying (assuming $C_{\ref{lem:sup-ratio}}$ is large enough so $s\ge 2$), and the first case follows from plugging in $s = kt/(24\pi)$ which satisfies $2 \le s \le k/(10\log{k})$ by the restriction on $t$ and $\delta$.
%If $48\pi/k\le t\le 1/\log k$, the infimum is achieved at $s = kt/(24\pi)$ and we obtain
%\[\mb{P}_{\vec{x}}\bigg[\sup_{\vec{a}\in A}\frac{\mb{P}_{\vec{b}}[\vec{b} = \vec{x}-\vec{a}]}{\mb{P}_{\vec{b}}[\vec{b} = \vec{x}]}\ge e^{tn}\bigg]\le |A|\exp\bigg(-\frac{kt^2n}{48\pi}\bigg).\]
%Now let $|A|\le\exp(\delta n)$, where $392\pi/k\le\delta\le\log 2$ (as clearly the statement allows such an alteration of $\delta$). 
Finally, since
\[0\le \sup_{\vec{a}\in A}\frac{\mb{P}_{\vec{b}}[\vec{b} = \vec{x}-\vec{a}]}{\mb{P}_{\vec{b}}[\vec{b} = \vec{x}]} \le \left(\max_{\ell}\frac{\max\Big\{\binom{k}{\ell- 1},\binom{k}{\ell}\Big\}}{\binom{k}{\ell}}\right)^{n} \le k^{n},\]
we have
\begin{align*}
\mb{E}_{\vec{x}}\bigg[\sup_{\vec{a}\in A}\frac{\mb{P}_{\vec{b}}[\vec{b} = \vec{x}-\vec{a}]}{\mb{P}_{\vec{b}}[\vec{b} = \vec{x}]}\bigg] 
&= \int_{-\infty}^{\log k}\mb{P}\bigg[\sup_{\vec{a}\in A}\frac{\mb{P}_{\vec{b}}[\vec{b} = \vec{x}-\vec{a}]}{\mb{P}_{\vec{b}}[\vec{b} = \vec{x}]}\ge e^{tn}\bigg]ne^{tn}dt\\
&\le \int_{1/\log{k}}^{\log{k}}\cdot + \int_{\sqrt{96\pi \delta/k}}^{1/\log{k}}\cdot + \int_{-\infty}^{\sqrt{96\pi\delta/k}}ne^{tn}dt\\
&\le e^{\sqrt{96\pi\delta/k}n} + \int_{\sqrt{96\pi\delta/k}}^{1/\log k}|A|\exp\bigg(-\frac{kt^2n}{48\pi}\bigg)ne^{tn}dt\\
&\quad+ \int_{1/\log k}^{\log k}|A|\exp\bigg(-\frac{kn}{48\pi(\log k)^2}\bigg)ne^{tn}dt\\
&\le\exp\left(O(\sqrt{\delta/k})n\right) + \int_{\sqrt{96\pi\delta/k}}^{1/\log{k}} ne^{-tn}dt + 1\\
&\le\exp\left(O(\sqrt{\delta/k})n\right).\qedhere
\end{align*}
\end{proof}





%%% AUTHOR: optional acknowledgments here
\section*{Acknowledgments} %%  you may comment this out if no Ackno
The authors are grateful to the anonymous reviewers for several suggestions which helped improve the presentation of the paper.
\providecommand{\bysame}{\leavevmode\hbox to3em{\hrulefill}\thinspace}
\providecommand{\MR}{\relax\ifhmode\unskip\space\fi MR }
% \MRhref is called by the amsart/book/proc definition of \MR.
\providecommand{\MRhref}[2]{%
  \href{http://www.ams.org/mathscinet-getitem?mr=#1}{#2}
}
\providecommand{\href}[2]{#2}
\begin{thebibliography}{1}

\bibitem{BHK09}
Andreas Bj{\"o}rklund, Thore Husfeldt, and Mikko Koivisto, \emph{Set
  partitioning via inclusion-exclusion}, SIAM Journal on Computing \textbf{39}
  (2009), no.~2, 546--563.

\bibitem{NPSW20}
Jesper Nederlof, Jakub Pawlewicz, Céline M.~F. Swennenhuis, and Karol
  Węgrzycki, \emph{A faster exponential time algorithm for bin packing with a
  constant number of bins via additive combinatorics}, pp.~1682--1701.

\bibitem{NV11}
Hoi Nguyen and Van Vu, \emph{Optimal inverse {L}ittlewood-{O}fford theorems},
  Adv. Math. \textbf{226} (2011), no.~6, 5298--5319. \MR{2775902}

\bibitem{RV08}
Mark Rudelson and Roman Vershynin, \emph{The {L}ittlewood--{O}fford problem and
  invertibility of random matrices}, Advances in Mathematics \textbf{218}
  (2008), no.~2, 600--633.

\bibitem{TV09}
Terence Tao and Van~H. Vu, \emph{Inverse {L}ittlewood--{O}fford theorems and
  the condition number of random discrete matrices}, Annals of Mathematics
  (2009), 595--632.

\bibitem{W17}
M{\aa}rten Wiman, \emph{Improved constructions of unbalanced uniquely decodable
  code pairs}, 2017,
  \url{https://www.diva-portal.org/smash/get/diva2:1120593/FULLTEXT01.pdf}.

\end{thebibliography}
\end{document}
